\chapter{Complexity — P}
%%% generated by the blueprint pipeline from TCSlib.Complexity.ClassP.P
%%%%%%%%%%%%%%%%%%%%%%%%%%%%%%%%%%%%%%%%%%%%%%%%%%%%%%%%%%%%%%%%%%%%%%%%%%%%%%%
\section{Overview}

This module defines the complexity class $\mathsf{P}$ of polynomial-time
decidable languages over the binary alphabet, following Arora--Barak
Definition~1.13, as the union over all degrees $c$ of the classes
$\mathrm{DTIME}(n^c + 1)$; the $+1$ padding keeps each component nonempty on the
empty input, where a bound of $n^c$ steps would vanish. The main results give
membership criteria for $\mathsf{P}$ from concrete polynomial bounds and certify
that the padded definition coincides with the conventional reading: a language
is in $\mathsf{P}$ exactly when some machine decides it within $C \cdot (n+1)^d$
steps for constants $C$ and $d$.

%%%%%%%%%%%%%%%%%%%%%%%%%%%%%%%%%%%%%%%%%%%%%%%%%%%%%%%%%%%%%%%%%%%%%%%%%%%%%%%
\section{Declarations}

\begin{definition}[The class of polynomial-time languages]
\lean{Complexity.P}
\label{Complexity.P}
\leanok
\uses{Complexity.DTIME}
The class $\mathsf{P}$ of polynomial-time decidable languages over the binary
alphabet: a language $L \subseteq \{0,1\}^{*}$ belongs to $\mathsf{P}$ when
there is a natural number $c$ such that $L \in \mathrm{DTIME}(n \mapsto n^c + 1)$,
i.e.\ some deterministic finite Turing machine decides $L$ within
$a \cdot (n^c + 1)$ steps on every input of length $n$, for some constant $a$.
In symbols, $\mathsf{P} = \bigcup_{c \in \mathbb{N}} \mathrm{DTIME}(n^c + 1)$.
\end{definition}

\begin{theorem}[Fixed-degree polynomial time is in P]
\lean{Complexity.dtime_poly_subset_P}
\label{Complexity.dtime_poly_subset_P}
\leanok
\uses{Complexity.DTIME, Complexity.P}
\difficulty{1}
For every natural number $c$, the class $\mathrm{DTIME}(n \mapsto n^c + 1)$ —
the languages over $\{0,1\}$ decided by some deterministic finite Turing machine
within $a \cdot (n^c + 1)$ steps on inputs of length $n$, for some constant $a$ —
is contained in $\mathsf{P}$.
\end{theorem}

\begin{theorem}[Membership in P from a pointwise polynomial bound]
\lean{Complexity.mem_P_of_dtime_le}
\label{Complexity.mem_P_of_dtime_le}
\leanok
\uses{Complexity.DTIME, Complexity.P, Complexity.dtime_poly_subset_P, Turing.FinTM.ComputesInTime.mono}
\difficulty{3}
Let $L$ be a language over the binary alphabet, let $T : \mathbb{N} \to \mathbb{N}$
be a time bound with $L \in \mathrm{DTIME}(T)$, and let $c$ and $d$ be natural
numbers such that $T(n) \le c \cdot (n^d + 1)$ for every $n$. Then
$L \in \mathsf{P}$. (The domination is required pointwise, at every input
length, not merely for all sufficiently large $n$.)
\end{theorem}

\begin{lemma}[Successor-power bound]
\lean{Complexity.succ_pow_le}
\label{Complexity.succ_pow_le}
\leanok
\difficulty{4}
For all natural numbers $n$ and $d$,
\[
(n + 1)^d \;\le\; 2^d \cdot \bigl(n^d + 1\bigr).
\]
\end{lemma}

\begin{theorem}[Characterization of P by conventional polynomial bounds]
\lean{Complexity.mem_P_iff}
\label{Complexity.mem_P_iff}
\leanok
\uses{Complexity.P, Complexity.succ_pow_le, Turing.FinTM, Turing.FinTM.ComputesInTime.mono, Turing.FinTM.DecidesInTime}
\difficulty{4}
A language $L$ over the binary alphabet belongs to $\mathsf{P}$ if and only if
there exist natural numbers $C$ and $d$ and a deterministic finite Turing
machine $M$ over $\{0,1\}$ such that $M$ decides $L$ within
$C \cdot (n + 1)^d$ steps on every input of length $n$. In particular the $+1$
padding in the definition of $\mathsf{P}$ describes the same class as the
conventional polynomial time bounds $C \cdot (n+1)^d$.
\end{theorem}

\begin{theorem}[Constant time is polynomial time]
\lean{Complexity.dtime_one_subset_P}
\label{Complexity.dtime_one_subset_P}
\leanok
\uses{Complexity.DTIME, Complexity.P, Complexity.mem_P_of_dtime_le}
\difficulty{1}
Every language over the binary alphabet in $\mathrm{DTIME}(n \mapsto 1)$ — that
is, decided by some deterministic finite Turing machine within a constant number
of steps, uniformly over all inputs — belongs to $\mathsf{P}$.
\end{theorem}
